\chapter{Endangered Species and Conservation}

After all, in the course of evolution, countless species have appeared and
disappeared so that the extinction of a single species cannot be the tragedy that some
would portray. Yet there is an atmosphere of sadness when a pretty plant such as
Primula kisoana nears extinction in the wild. It is here that the amateur and
professional grower can make a contribution. The grower can accumulate every clone
available and can increase the genetic pool and evolutionary health of the species by
a program of cross pollination, propagation, and distribution. In part this book is
dedicated to such programs.

Such propagation programs compliment programs directed towards preserving
natural habitats. The Nature Conservancy has done a magnificent job in studying the
needs of species, both plants and animals, and setting up sanctuaries by land
purchases and other means. But this is not enough. Even now a number of species
are more plentiful in cultivation than in the wild; Gingko, Metasequoia, and Primula
kisoana for examples. With the present density of human population in the world it is
futile to think that nature preserves of such magnitude can be established that would
insure the survival of the above three species by natural propagation. The possibility
that they are naturally dieing out cannot be overlooked.

The United States government and State governments have published lists of
endangered species. It is illegal to distribute these plants or seeds of these plants.
This policy and the laws designed to enforce this policy can only serve to
\textbf{speed up the extinction of the species by discouraging propagation}. Further, many of
the species on the present list are only minor variants of common species. The law
can be easily circumvented by distributing the variant under the name of the species,
particularly since the variants are marginally distinct, or by distributing under a new
name since there is no final authority at present on names of species. It is felt that the
laws should be more specifically directed towards preventing commercial collecting.

The laws could also be rewritten to encourage propagation, especially by seed.

Two groups oppose propagation programs. One group places their force
behind saving plants as if plants were static objects. In their view it becomes more
important to save plants than to grow plants. In other words a plant saved is better
than a plant grown. A second group has the view that raising large numbers from
seed, introducing the species into new and suitable habitats, and other actions which
promote and accelerate evolution are looked upon at best as causing taxonomic grief.
This book takes the view that these oppositions to propagation are in error. Our
society already propagates, manages, and distributes animals and fish with much
success. Now we must manage plants the same way. This is the way to preserve
endangered species and serve conservation, and most of all, to express our love of
plants.

